\clearpage
\section*{September 9, 2026: Peoria Green and individual townhome lots}
\addcontentsline{toc}{section}{September 9, 2026: Peoria Green and individual townhome lots}
\textit{Agent-drafted entry.}

Peoria Green is retained as one two-wing project with 17 homes completed in 2021, following Jacob's approval. Completed City revision permit 100876378 reduces the original 20-home proposal to 17. Contemporary completion reporting identifies 12 homes at 128 S Green and five at 123 S Peoria. All 17 entries in the completed 2022 Assessor condominium base repeat 56,295 square feet of building area and 16,695 square feet of land; these totals enter once. Four parking flags had reduced the counted homes to 13. At least one flagged entry reports five bedrooms and 3,859 square feet, contradicting a parking-only interpretation.

A general rule resolves seven townhomes at 1911--1923 S Calumet and the four-home building at 1705 W Chicago. Their older tax records referenced their own property numbers as if they required combination with another property. Later independent assessments contain one complete building record, with unchanged property number, class, construction year, floor area and residential count. The rule requires a one-property historical lineage and no concurrent multiple cards. It uses the whole later report, preferring 2025, rather than combining measurements from different assessments or calculating land from tax shares.

The seven Calumet homes retain 2009 and their reported floor areas: 2,584 square feet each for 1911--1921 and 2,790 for 1923. The later Assessor reports six individual lots of 1,000 square feet and one of 1,500. The frozen paper file instead repeats 51,456--224,142 square feet of development land on these individual homes, understating density. These seven frozen observations are outside 500 feet of a ward boundary. Peoria's frozen record is about 96 feet away and counts 13 homes; its eventual correction to 17 can therefore affect the 500-foot sample. No regression effect has been estimated.

Two other cases follow existing exclusion rules. Completed permit 100757281 describes Tribune Tower at 435 N Michigan as adaptive reuse of an existing office/retail building into 162 condos, with an addition; it is not ground-up construction. At 1920 W Cullerton, the uniquely matched finished condominium base contains three residential entries, all reporting 2023. The completed-Assessor-year rule holds it outside 2006--2022 even though its earlier residential record says 2022.

The fixed 268-record review now has 185 closed questions and 83 open, down from 94 open. These eleven resolutions add nine selected residential projects, bringing the file from 12,909 to 12,918, and explicitly exclude two others. All nine retained projects have construction-year boundary distances and no commercial component-PIN overlap. Previously selected residential measurements and all 815 selected commercial measurements are unchanged. The construction task and targeted reconciliation report build passed; an unchanged construction build runs no substantive producer.

\paragraph{Sources and reproduction.} The existing Assessor histories and City permit snapshot supply the numerical inputs. External corroboration is recorded in the decision ledger: the August 28, 2021 Chicago YIMBY Peoria completion article and Golub's Tribune Tower project description. Changes remain on \texttt{sales\_sample\_exploration}. The complete Peoria decision replaces its earlier spatial-link-only exception. The paper's frozen input remains unchanged; the combined dataset and remaining reviews are unfinished.
